\documentclass[11pt,a4paper]{article}
\usepackage[T1]{fontenc}
\usepackage{lmodern,amsmath,amssymb,amsthm}
\usepackage[margin=28mm]{geometry}
\newcommand{\E}{\mathbb E}
\newcommand{\ind}{\mathbf1}
\newcommand{\R}{\mathbb R}
\begin{document}
\section*{Analytic soundness of the optional coefficient audit}
Let $\E X=0$, $\E X^2=1$, with $-R_-\le X\le R_+$.
Write $\rho=\E|X|^3$, $\gamma=\E X^3$, $\beta=\E[X|X|]$,
$K_1=\E[X|X|^3]-3\beta-3\rho\gamma/2$.
For the actual sign split let $V=\E[X\mid\ind_{\{X>0\}}]$,
$W=X-V$ and $\theta_\pm=\E[W^2;\pm X>0]$.
The signed conditional and central-gap clauses below assume an already
proved gap $(-A,B)$ with $A,B>0$; the moment matrices also allow an atom at zero.
Conditional centering proves, for signwise constant $h$ and $j\ge2$,
\[
 \E[h(|X|^j-|V|^j)]
 =\E\left[hW^2\sum_{k=0}^{j-2}(k+1)|V|^k|X|^{j-2-k}\right].
\]
Indeed multiply the quotient by $(|X|-|V|)^2$ and restore
$j|V|^{j-1}(|X|-|V|)$, whose conditional expectation is zero.
For every support gap $(u,v)$ put $\tau_{u,v}=\E[X;X>t]$ at
$t\in(u,v)$, and define
\[
 J_1=\sum\tau_{u,v}\int_u^v(t-u)t\,dt,\qquad
 J_2=\sum\tau_{u,v}\int_u^v(t-u)t|t|\,dt.
\]
Use subscripts $c,o$ for the central and outer gap contributions.
Let $m=\E X_+=\E X_-$,
$T_+=\E X_+^3-mB^2$ and $T_-=\E X_-^3-mA^2$.
The common nonnegative outer budget $\mathcal B$ appearing below is
an explicit input hypothesis, not acquired by polynomial acceptance.
Here $\sigma=\sqrt{10}$; any proved positive lower prices may replace
$\sigma-1,\sigma+1$ in a weaker budget.

% Independent analytic license for the algebraic coefficient family.
% The surrounding text supplies the coefficient contract and existing notation.
\paragraph{Why the coefficient formulas bound one marginal.}
There are three underlying calculations. They use the same signwise
moments, conditional variances and gap budget throughout; the exact
coefficient object records their polynomial expansions.

First, on either sign let $Y=|X|\in[0,r]$ and write
$s_j=\E[Y^j;\text{that sign}]$. Both signs have $s_1=m$.
For $H_j=\left(\begin{smallmatrix}s_j&s_{j+1}\\s_{j+1}&s_{j+2}\end{smallmatrix}\right)$,
\[
 H_j\succeq0\quad(0\le j\le3),\qquad
 rH_j-H_{j+1}\succeq0\quad(0\le j\le2).
\]
Their quadratic forms integrate, respectively,
$Y^j(a+bY)^2$ and $Y^j(r-Y)(a+bY)^2$. These identities prove the
whole moment and support-contraction family, separately at the two
support radii. A quotient deduction from
$\left(\begin{smallmatrix}A&B\\B&C\end{smallmatrix}\right)\succeq0$
uses $C\ge B^2/A$ only when $A>0$. At $A=0$ it uses $B=0$, $C\ge0$;
an interval meeting zero may simply discard the nonnegative correction.
It may not divide by zero or square an unlicensed signed endpoint.

Second, the conditional Taylor identity
displayed above integrates its polynomial remainder
against the nonnegative measures
$d\nu_\pm=(X-V)^2\ind_\pm dP$, of masses $\theta_\pm$.
With $\bar a_\pm=|\E[X\mid\pm]|$, $\gamma_V=\E V^3$, it gives
\begin{gather*}
 K_1(X)-K_1(V)=
 \int F(x,\bar a_+;\rho+\gamma_V)\,d\nu_+
 -\int F(x,\bar a_-;\rho-\gamma_V)\,d\nu_-,\\
 F(x,a;t)=x^2+2ax+3a^2-3-\tfrac32t(x+2a).
\end{gather*}
Indeed expand the product as
$\rho\gamma-\rho_V\gamma_V=\rho(\gamma-\gamma_V)
+\gamma_V(\rho-\rho_V)$ and use the conditional identities at orders
two, three and four. The factored derivatives certify $F_t\le0$ and
the positive Hessian $2(dx+da)^2+4(da)^2$.
Thus a rectangle maximum is at a corner; a minimum is among its corners,
its clipped edge stationary points and its interior stationary point.
All are included in the exact coefficient map. If the two signed
integrands lie in $[L,U]$, their sum lies in
$[\Theta\min(0,L),\Theta\max(0,U)]$ when
$\theta_++\theta_-\le\Theta$. This includes zero noise without division.

Third, direct integration over the central gap $(-A,B)$ gives
\begin{gather*}
 J_{1,c}=mP_1(A,B),\quad J_{2,c}=mP_2(A,B),\\
 P_1=B^3/3+AB^2/2-A^3/6,\qquad
 P_2=B^4/4+AB^3/3-A^4/12.
\end{gather*}
Both polynomials increase in $B$ and their $A$ derivatives change sign
only at $A=B$. Their extrema therefore use the indicated endpoints and
the clipped point $A=B$. For $G=\rho P_1-P_2$, the complete central
bound uses the degree-four tensor Bernstein coefficients on the closed
$(A,B)$ rectangle. The basis is nonnegative and sums to one, so retain
both $m_l$ and $m_u$ times every coefficient. This handles negative
coefficients and zero-width rectangles. The result is a maximum of
affine functions of $\rho$, hence convex. The checked inverse basis
identities prove coverage of the entire rectangle, not only a grid.

For a positive outer gap $(a,b)$, $a\ge B$, the square inequality
$(t-a)(\rho-t)\le(\rho-B)_+^2/4$ and signwise Tonelli give
\[
 3(\rho J_1-J_2)_o\le\tfrac38
 \{(\rho-B)_+^2T_++(R_--\rho)_+^2T_-\},\qquad
 -R_-T_-/2\le J_{1,o}\le R_+T_+/2.
\]
For the first bound on a negative gap $(u,v)$, put $x=-t$ and $b=-u$:
its integrand is $(b-x)x(x-\rho)$, whose positive part is at most
$x(R_--\rho)_+^2/4$. Enlarge each union of gaps to its complete
outer ray; the integrand is nonnegative, and Tonelli gives
$\int_B^\infty t\tau(t)\,dt=T_+/2$ and its negative counterpart.
This proves the factor $3/8$ even if the support has positive Lebesgue
measure. For the second bound, $0\le t-u\le R_+$ on positive outer
gaps and $0\le t-u\le R_-$ on negative outer gaps; integrate
$|t|\tau(t)$ over the complete corresponding ray.
Endpoint atoms do not change these Lebesgue integrals.

These costs share the one budget
$(\sigma-1)T_++(\sigma+1)T_-\le \mathcal B$ as a separate premise. Intersect it with any proved individual
caps before pricing a linear cost. Without caps, use
\[
 b_+T_++b_-T_-\le(\mathcal B)
 \max\{0,b_+/(\sigma-1),b_-/(\sigma+1)\}.
\]
The zero option is necessary for signed costs. Sum costs before this
single pricing step; intersect alternative upper bounds afterwards.
Every coefficient interval is a necessary condition on this one law.
The analytic formation of the central gap and the prior support bounds
remain prerequisites; coefficient acceptance does not establish them.

\end{document}
